\section*{Hoofdstuk \ref{ch:b2:meiosis-heredity} --- Meiose, genetische menging en erfelijkheid}

\begin{solution}{exo:b2:meiosis-heredity:1}
Mitose: één deling, geen paring van homologen, twee cellen, elk
genetisch identiek aan de oudercel ($2n \to 2n$). \omterm{def:b2:meiosis-heredity:meiosis}{Meiose}: twee delingen
na één replicatie, homologen paren en ondergaan crossing-over, vier
cellen, \omterm{def:b2:meiosis-heredity:meiosis}{haploïd} en genetisch alle verschillend ($2n \to n$).
\end{solution}

\begin{solution}{exo:b2:meiosis-heredity:2}
$2^{23} = 8.4\times 10^{6}$; $2^{4} = 16$; $2^{7} = 128$ --- nog voordat
crossing-over elk van deze aantallen onbeperkt vermenigvuldigt.
\end{solution}

\begin{solution}{exo:b2:meiosis-heredity:3}
$Tt\times Tt$: \omterm{def:b2:meiosis-heredity:vocabulary}{genotypen} $1\,TT : 2\,Tt : 1\,tt$, \omterm{def:b2:meiosis-heredity:vocabulary}{fenotypen} $3$ hoog
: $1$ dwerg. $Tt\times tt$: $1\,Tt : 1\,tt$, $1 : 1$. Kruis de hoge
plant terug met een dwergplant: allemaal hoog betekent $TT$, de helft
dwerg betekent $Tt$ (of laat haar zichzelf bestuiven: $TT$ is raszuiver,
$Tt$ splitst uit in $3 : 1$).
\end{solution}

\begin{solution}{exo:b2:meiosis-heredity:4}
Gekoppelde genen liggen op hetzelfde chromosoom en hun parentale
allelcombinaties zijn onder de gameten oververtegenwoordigd; de
\omterm{def:b2:meiosis-heredity:linkage}{recombinatiefrequentie} is de fractie recombinante gameten (recombinante
nakomelingen van een testkruising); één \omterm{def:b2:meiosis-heredity:linkage}{centimorgan} is één procent
recombinatie. Genen die ver uit elkaar liggen op een lang chromosoom
hebben in vrijwel elke \omterm{def:b2:meiosis-heredity:meiosis}{meiose} minstens één crossing-over tussen zich, en
meervoudige crossing-overs maken de combinaties willekeurig, zodat $r$
de waarde $\tfrac12$ bereikt --- niet te onderscheiden van \omterm{thm:b2:meiosis-heredity:mixing}{onafhankelijke
verdeling}.
\end{solution}

\begin{solution}{exo:b2:meiosis-heredity:5}
Totaal 556; verwacht $312.75 : 104.25 : 104.25 : 34.75$; $\chi^2 =
2.25^2/312.75 + 3.75^2/104.25 + 3.25^2/104.25 + 2.75^2/34.75 = 0.016
+ 0.135 + 0.101 + 0.218 = 0.47 < 7.81$: verenigbaar met $9 : 3 : 3 :
1$.
\end{solution}

\begin{solution}{exo:b2:meiosis-heredity:6}
Tegen $1 : 1 : 1 : 1$ (elk 250): $\chi^2 = (162^2 + 138^2 + 147^2
+ 153^2)/250 = 361 \gg 7.81$: gekoppeld. Recombinanten $103 + 97 = 200$
op 1000: $r = 0.20$, \qty{20}{cM}. In de afstotingsfase geeft $Ab/aB$ 400
$Ab$, 400 $aB$, 100 $AB$, 100 $ab$: dezelfde afstand, met de parentale
en recombinante klassen verwisseld.
\end{solution}

\begin{solution}{exo:b2:meiosis-heredity:7}
$d = -\tfrac12\ln(1 - 2r)$: $r = 0.20$ geeft $-\tfrac12\ln 0.6 =
0.255$, \qty{25.5}{cM}; $r = 0.45$ geeft $-\tfrac12\ln 0.1 = 1.15$,
\qty{115}{cM}. $r = \tfrac12(1 - e^{-2d})$: \qty{30}{cM} geeft
$\tfrac12(1 - e^{-0.6}) = 0.226$; \qty{100}{cM} geeft $\tfrac12(1 -
e^{-2}) = 0.43$. Onder \qty{10}{cM} is de correctie verwaarloosbaar;
boven \qty{30}{cM} is ze onmisbaar, en boven \qty{50}{cM} raakt $r$
verzadigd, zodat één tweepuntskruising de afstand niet meer kan meten
--- verre genen brengt men in kaart door nabije aan elkaar te rijgen.
\end{solution}

\begin{solution}{exo:b2:meiosis-heredity:8}
(a) $X^wX^w \times X^+Y$: alle zonen wit ($X^wY$), alle dochters rood
($X^+X^w$). (b) $X^+X^w \times X^wY$: zonen half rood, half wit;
dochters half rood ($X^+X^w$), half wit ($X^wX^w$). (c) dochters
$X^+X^w$ $\times$ broers $X^wY$: zoals in (b) --- de helft van elk geslacht
wit, de helft rood.
\end{solution}

\begin{solution}{exo:b2:meiosis-heredity:9}
Twee enzymen in serie maken het pigment; paars vereist een \omterm{def:b2:meiosis-heredity:vocabulary}{dominant}
\omterm{def:b2:meiosis-heredity:vocabulary}{allel} op beide loci ($C\_P\_$). De lijnen zijn $CCpp$ en $ccPP$
(elk wit om een andere reden); de hybride $CcPp$ is paars;
na zelfbestuiving geeft hij $9\,C\_P\_$ paars op $7$ wit ($3\,C\_pp + 3\,ccP\_ +
1\,ccpp$). Testkruising $CcPp\times ccpp$: $1$ paars ($CcPp$) : $3$
wit.
\end{solution}

\begin{solution}{exo:b2:meiosis-heredity:10}
Parentaal $+++$, $abc$; dubbele crossing-overs $+b+$, $a+c$; $b$ is het
gen dat wisselde: volgorde $a$--$b$--$c$. $r_{ab} = (11 + 9 + 2)/1202
= 0.018$ (\qty{1.8}{cM}); $r_{bc} = (118 + 112 + 2)/1202 = 0.193$
(\qty{19.3}{cM}). Verwachte dubbele $0.018\times 0.193\times 1202 =
4.2$; waargenomen 2: coïncidentie $0.47$, interferentie $0.53$.
\end{solution}

\begin{solution}{exo:b2:meiosis-heredity:11}
Asci met segregatie bij de tweede deling $300/1000 = \qty{30}{\%}$; afstand $=
\tfrac12\times 30 = \qty{15}{cM}$. Een crossing-over tussen gen en
centromeer betreft maar twee van de vier chromatiden, zodat een ascus
met segregatie bij de tweede deling twee recombinante en twee
parentale chromatiden draagt: de \omterm{def:b2:meiosis-heredity:linkage}{recombinatiefrequentie} is de helft van
de frequentie van zulke asci. Voor $2 : 4 : 2$: een crossing-over tussen
de locus en het centromeer verwisselt de twee binnenste chromatiden, zodat
elke halve spoelfiguur in \omterm{def:b2:meiosis-heredity:meiosis}{meiose} I één chromatide $A$ en één chromatide
$a$ draagt; \omterm{def:b2:meiosis-heredity:meiosis}{meiose} II scheidt ze dan in de volgorde $A\,a \mid
a\,A$ (of omgekeerd), en de laatste mitose verdubbelt elke spore:
$2 : 4 : 2$.
\end{solution}

\begin{solution}{exo:b2:meiosis-heredity:12}
Haar moeder is een obligate draagster (zij kreeg de enige X van haar
vader), zodat de vrouw het \omterm{def:b2:meiosis-heredity:vocabulary}{allel} met kans $\tfrac12$ heeft ontvangen.
Twee gezonde zonen: een draagster heeft kans $(\tfrac12)^2 = \tfrac14$
op twee gezonde zonen, een niet-draagster kans 1. Bayes:
$P(\text{draagster}\mid\text{gegevens}) = (\tfrac12\times\tfrac14)/(\tfrac12
\times\tfrac14 + \tfrac12\times 1) = \tfrac{1/8}{5/8} = \tfrac15$.
\end{solution}

\begin{solution}{pb:b2:meiosis-heredity:1}
\textbf{1.} Ouders $vg/vg\;+/+$ en $+/+\;e/e$; $F_1$ $+/vg\;+/e$;
gameten $+\,+$, $+\,e$, $vg\,+$, $vg\,e$, elk $\tfrac14$.
\textbf{2.} $900 : 300 : 300 : 100$.
\textbf{3.} $\chi^2 = 8^2/900 + 9^2/300 + 4^2/300 + 3^2/100 = 0.07 +
0.27 + 0.05 + 0.09 = 0.48 < 7.81$: de \omterm{thm:b2:meiosis-heredity:mixing}{onafhankelijke verdeling} klopt.
\textbf{4.} \omterm{def:b2:meiosis-heredity:vocabulary}{Homozygoot} wild op beide loci is $\tfrac{1}{16}$ van het geheel,
$\tfrac{9}{16}$ is wild: $\tfrac19$.
\textbf{5.} Wild, rudimentair, ebben, rudimentair ebben, elk $\tfrac14$.
\textbf{6.} Het \omterm{def:b2:meiosis-heredity:vocabulary}{fenotype} is het product van \omterm{def:b2:meiosis-heredity:vocabulary}{genotype} en omgeving
(de temperatuur werkt in op de pigmentenzymen): een \omterm{def:b2:meiosis-heredity:vocabulary}{genotype} legt een
reactienorm vast, geen vast kenmerk.
\textbf{7.} $b\,+/+\,vg$ (afstotingsfase): $b$ kwam met $+_{vg}$ van de ene
ouder, $vg$ met $+_b$ van de andere.
\textbf{8.} Parentaal: zwart ($b\,+$) 965 en rudimentair ($+\,vg$) 944;
recombinant: zwart rudimentair ($b\,vg$) 206 en wild ($+\,+$) 185.
Het wildtypechromosoom $+\,+$ bestond niet in de $F_1$; het kan
alleen door een crossing-over worden gemaakt.
\textbf{9.} $r = 391/2300 = 0.17$: \qty{17}{cM}.
\textbf{10.} $d = -\tfrac12\ln(1 - 0.34) = -\tfrac12\ln 0.66 = 0.21$:
\qty{21}{cM}.
\textbf{11.} Bij mannelijke \emph{Drosophila} is er geen crossing-over:
de gameten van het mannetje dragen alleen de twee parentale chromosomen.
\textbf{12.} $F_1$ $b\,vg/+\,+$ (koppelingsfase): parentale klassen wild en
zwart rudimentair, elk \qty{41.5}{\%}; recombinant zwart en
rudimentair, elk \qty{8.5}{\%}.
\textbf{13.} $X^hX^+$: een dochter ontvangt de enige X van haar vader, die
$h$ draagt, en een normale X van haar niet-aangedane moeder.
\textbf{14.} $\tfrac12$ (zoon) $\times$ $\tfrac12$ (ontvangt $X^h$) $=
\tfrac14$.
\textbf{15.} $\tfrac12$: ze ontvangt een van beide X-chromosomen van de
moeder; de vader geeft een normale X.
\textbf{16.} Elk kind is niet-aangedaan met kans $\tfrac34$:
$(\tfrac34)^3 = \tfrac{27}{64} = 0.42$.
\textbf{17.} Een zoon ontvangt de Y van zijn vader, nooit diens X. Een vrouw
is alleen aangedaan als $X^hX^h$: een aangedane vader en een draagster (of
aangedane) moeder, wat zeldzaam is.
\textbf{18.} A-prioriwaarde $\tfrac12$. Twee gezonde zonen hebben kans
$\tfrac14$ als ze draagster is, 1 als ze dat niet is: a posteriori $(\tfrac12\times
\tfrac14)/(\tfrac12\times\tfrac14 + \tfrac12) = \tfrac15$.
\textbf{19.} De enige X van een zoon komt van zijn moeder, en de Y draagt
geen van deze genen, zodat elke zoon één maternale gameet ongemaskeerd
toont: de kruising is haar eigen testkruising.
\textbf{20.} Parentaal $+++$ (853), $ywm$ (833); dubbele crossing-overs
$+w+$ (2), $y+m$ (2); $w$ wisselde: volgorde $y$--$w$--$m$.
\textbf{21.} $r_{yw} = (24 + 26 + 2 + 2)/2000 = 0.027$; $r_{wm} = (128
+ 132 + 2 + 2)/2000 = 0.132$. Kaart: $y$ --- \qty{2.7}{cM} --- $w$ ---
\qty{13.2}{cM} --- $m$.
\textbf{22.} Verwacht $0.027\times 0.132\times 2000 = 7.1$; waargenomen
4: coïncidentie $0.56$, interferentie $0.44$.
\textbf{23.} Recombinanten $y$--$m$: $y++$, $+wm$, $yw+$, $++m$: $(24 +
26 + 128 + 132)/2000 = 0.155$, tegen $0.027 + 0.132 = 0.159$: de
vier dubbele crossing-overs herstellen de parentale combinatie $y$--$m$ en
worden gemist, $2\times 4/2000 = 0.004$.
\textbf{24.} Dochters ontvangen de X $y\,w\,m$ van de vader en één
maternale X: ze zijn allemaal \omterm{def:b2:meiosis-heredity:vocabulary}{heterozygoot} en fenotypisch wild, welke
recombinante X ze ook dragen; de recombinatie is onzichtbaar.
\textbf{25.} $y$ --- \qty{2.7}{cM} --- $w$ --- \qty{13.2}{cM} --- $m$;
coïncidentiecoëfficiënt $0.56$, interferentie $0.44$.
\end{solution}
